% Generated draft: requires mathematical and source review.
\documentclass[11pt]{article}
\usepackage[margin=25mm]{geometry}
\usepackage{fontspec}
\setmainfont{Latin Modern Roman}
\usepackage{amsmath,amssymb,mathtools,unicode-math}
\setmathfont{Latin Modern Math}
\usepackage{enumitem,needspace,xurl,hyperref}
\setlength{\parindent}{0pt}
\setlength{\parskip}{5pt}
\setlength{\emergencystretch}{3em}
\newcommand{\N}{\mathbb{N}}
\newcommand{\Z}{\mathbb{Z}}
\newcommand{\Q}{\mathbb{Q}}
\newcommand{\R}{\mathbb{R}}
\newcommand{\C}{\mathbb{C}}
\newcommand{\F}{\mathbb{F}}
\newcommand{\A}{\mathbb{A}}
\newcommand{\PP}{\mathbb{P}}
\newcommand{\E}{\mathbb{E}}
\newcommand{\eps}{\varepsilon}
\newcommand{\dd}{\,\mathrm{d}}
\begin{document}
\section*{UnsolvedMath ID: 6700001}
\textbf{Source problem code:} catalogue entry
\subsection{Definitions and mathematical statement}
For a smooth manifold \(M\) and a smooth Riemannian metric \(g\), let \(\operatorname{Scal}_g:M\to\R\) be the trace of the Ricci tensor with respect to \(g\). Say that \(g\) has positive scalar curvature when
\[
\operatorname{Scal}_g(x)>0\quad\text{for every }x\in M.
\]
Gromov asks: which topological types of manifolds admit such a metric? For any specified class \(\mathcal C\) of smooth manifolds, this asks to characterize
\[
\{[M]_{\mathrm{homeomorphism}}: M\in\mathcal C,\ \exists g\text{ Riemannian on }M\text{ with }\operatorname{Scal}_g>0\}.
\]
The source does not specify \(\mathcal C\): dimension, compactness, absence of boundary, orientability, and completeness of the metric are not stated. Nor does it specify whether ``topologies'' means homeomorphism types, diffeomorphism types, or another classification. The display gives one precise interpretation after choosing those conventions; it does not restrict the question to closed or simply connected manifolds. A topological manifold also needs a smooth structure before the metric condition can be imposed, and existence for some smooth structure must be distinguished from existence for a fixed one.
\subsection{Short English statement}
Classify the manifolds that support a metric with positive scalar curvature at every point. The intended class of manifolds and the meaning of topological classification need clarification.
\subsection{Sources}
\begin{itemize}
\item Local UnsolvedMath registry, permanent ID 6700001, source code catalogue entry; attributed to Misha Gromov, 2017; updated 2026-08-06.
\item Gromov, \emph{101 Questions, Problems and Conjectures around Scalar Curvature} (2017), printed tag [?1], occurrence 1, PDF page 1: \url{https://www.ihes.fr/~gromov/wp-content/uploads/2018/08/101-problemsOct1-2017.pdf}. The registry records access on 2026-07-29 and normalized PDF text extraction. It identifies the PDF as public, author-hosted, and incomplete, with present status and reuse terms needing review.
\item The supplied research summary cites Schoen--Yau (1979), Gromov--Lawson (1980), Perelman/Ricci flow, and Stolz (1992, \emph{Annals of Mathematics}), giving the latter title as ``A conjecture concerning positive Ricci curvature'' while describing a scalar curvature index classification. These are supplied attributions with incomplete or potentially inconsistent bibliographic details, not independently checked citations.
\end{itemize}
\end{document}
